\documentclass{article}
\usepackage{amsmath,amssymb}
\usepackage{xurl}
\usepackage[colorlinks=true,linkcolor=blue,citecolor=blue,urlcolor=blue]{hyperref}
% Shared commands for notes in docs/paper-gaps/.

\DeclareUnicodeCharacter{2080}{\ifmmode{}_{0}\else\textsubscript{0}\fi}
\DeclareUnicodeCharacter{2081}{\ifmmode{}_{1}\else\textsubscript{1}\fi}
\DeclareUnicodeCharacter{2082}{\ifmmode{}_{2}\else\textsubscript{2}\fi}
\DeclareUnicodeCharacter{2099}{\ifmmode{}_{n}\else\textsubscript{n}\fi}

\newcommand{\Lean}{\textsc{Lean}}
\ExplSyntaxOn
\NewDocumentCommand{\leanid}{m}{
  \ifmmode
    \textsf{\detokenize{#1}}
  \else
    \group_begin:
    \urlstyle{sf}
    \tl_set:Nn \l_tmpa_tl {#1}
    \tl_replace_all:Nnn \l_tmpa_tl {\_} {_}
    \exp_args:NV \path \l_tmpa_tl
    \group_end:
  \fi
}
\ExplSyntaxOff
\providecommand{\tr}{\operatorname{tr}}
\newcommand{\tnleanIssueBase}{https://github.com/LionSR/TNLean/issues/}
\newcommand{\tnleanPrBase}{https://github.com/LionSR/TNLean/pull/}
\newcommand{\tnleanDocBase}{https://sirui-lu.com/TNLean/docs/}
\newcommand{\ghissue}[1]{\href{\tnleanIssueBase#1}{GitHub issue \##1}}
\newcommand{\ghpr}[1]{\href{\tnleanPrBase#1}{PR \##1}}
\newcommand{\doclink}[1]{\href{\tnleanDocBase#1.html}{\path{#1}}}

% Dirac notation and the identity operator, as in blueprint/src/macros/common.tex.
% \providecommand keeps note-local definitions authoritative.
\usepackage{dsfont}
\providecommand{\ket}[1]{|#1\rangle}
\providecommand{\bra}[1]{\langle #1|}
\providecommand{\braket}[2]{\langle #1 | #2 \rangle}
\providecommand{\ketbra}[2]{|#1\rangle\!\langle #2|}
\providecommand{\Id}{\mathds{1}}
\providecommand{\one}{\mathds{1}}
\providecommand{\C}{\mathbb{C}}
\providecommand{\MN}[1]{M_{#1}(\C)}
\providecommand{\MD}{M_D(\C)}

% Verdict marker: \gapnote{<kind>}{<status>}. Kind is the policy
% classification (clarification, local-correction, scope-restriction,
% unfaithful, false-source, open-gap); status is open, wip (elimination
% underway), resolved, or historical. Machine-read by the site index;
% prints a verdict line.
\newcommand{\gapnote}[2]{\par\noindent\textbf{Verdict:} #1 (#2).\par}

\providecommand{\gapnote}[2]{\par\noindent\textbf{Verdict:} #1 (#2).\par}
\title{Hastings' High-Temperature Product Form on Planar Domains}
\author{}
\date{October 23, 2026}
\begin{document}
\maketitle
\gapnote{scope-restriction}{open}

\section{The source setting}

Section~IV of Hastings' \emph{Solving gapped Hamiltonians
    locally}~\cite{Hastings2006Solving} (arXiv:cond-mat/0508554v3,
\texttt{proj3new.tex}, lines 1017--1079) works on an arbitrary lattice with a
metric $d$, interaction terms $H_i$ supported within distance $R$ of a site $i$,
and at most $S$ sites within distance $R$ of any site. The matrix product
form displayed at source lines 1024--1039 asks for coefficient functions $F_i$ that depend only on the
indices $\alpha_j$ with $d(i,j)\le l_{\mathrm{proj}}$. The source claims that
every Trotter product $\rho_n=[\prod_i e^{-\beta H_i/n}]^n$ has this form
exactly, with range $R$, and that after grouping the indices of all factors at
a site each site carries at most $nS$ indices (lines 1069--1075).

\section{The formalized setting}

The formalization states the product form for the local Hamiltonians already
used for the two-dimensional area law: the sites form a finite domain
$\Lambda\subseteq\mathbb Z^2$, the distance between two sites is the length of
a shortest walk joining them in the nearest-neighbour graph of $\Lambda$, and a
term of range $R$ is supported on a set of sites any two of which are joined by
a walk of length at most $R$. This is a diameter bound, whereas the source
assumes a radius bound about the specified centre of each term. A support
contained in a ball of radius $R$ can have diameter $2R$. Thus the current
statement is a specialization with a stricter support condition; it does not
formalize the source theorem with the same range parameter. Replacing the
support diameter by $2R$ would yield range $2R$ by the current argument,
not the source's range $R$.

\paragraph{Elimination plan.} First define the product form on an arbitrary
finite metric site set, with coefficient locality measured by its given
metric, independently of a graph or a planar embedding. Next prove the
supported-operator expansion with a specified centre $i$ and the sole
support condition $X\subseteq\{j:d(i,j)\le R\}$. Attach the nonconstant
matrix-unit coefficient to $i$, retaining the diagonal identity factors off
$X$ even when $i\notin X$. This removes the diameter assumption while
preserving the source's radius $R$. Finally prove product closure and the
Trotter theorem for terms indexed by their centres under these radius
hypotheses. Derive the present diameter-bounded statement by choosing any
site of each nonempty support as its centre and specializing the geometry.
The source's index-component count requires a separate construction that
assigns components only to sites touched by a factor and uses the bound $S$
on each radius-$R$ ball. These extensions are part of
\href{https://github.com/LionSR/TNLean/issues/9854}{issue 9854}.

\section{What is not claimed}

The formalization proves the existence of the product form with range $R$ for
every Trotter product whose term supports have diameter at most $R$. It does not establish the source's count of the indices
at a site. The source groups at each site one index component for every
occurrence of a factor acting near that site, at most $nS$ components. Here
each factor $e^{-\beta H_i/n}$ is expanded in single-site matrix units at every
site, so a Trotter product of $m$ factors on local dimension $q$ carries $m$
components at each site, one for every factor wherever it acts, each with
$q^2$ values. Any later use of the
product form that needs the index count must establish it separately.

\bibliographystyle{unsrt}
\bibliography{references}
\end{document}
